\chapter{Introduction}

Supernumerary robotic limbs (SRLs) are additional arms worn on the body, which
act alongside the wearer's own limbs rather than replacing or rehabilitating
them. They differ from prostheses in that the wearer retains full natural
function, and from industrial manipulators in that the base of the robot is a
person: it accelerates, sways and steps, and it is soft, valuable and
irreplaceable. The work reported in this thesis concerns a wearable system of
two seven-degree-of-freedom Kinova Gen3 arms carried on a backpack frame and
teleoperated from an instrumented mannequin master arm.

The configuration that this system implements is deliberately a
\emph{two-person} one: the arms are worn by one individual and commanded by
another. \todo{CITE: Fusion, Saraiji et al., SIGGRAPH 2018 Emerging
Technologies -- the direct predecessor for the operator/wearer architecture}
That separation is the source of most of the interesting engineering
difficulty in the project, because the person who bears the risk has no
control input, and the person issuing the commands has neither proprioception
of the wearer nor any vestibular sense of the base motion they are working
against.

\section{Objectives}

The project set out to build and characterise a complete teleoperation stack
for this configuration, and specifically to establish, by measurement rather
than by demonstration, what such a system can and cannot do. The objectives
were:

\begin{enumerate}
  \item to implement an end-to-end pipeline from the master arm's sensors to
    joint trajectories on two arms, including the mapping from the operator's
    limb pose to a commanded end-effector pose;
  \item to design a safety layer appropriate to a robot mounted beside a
    person's head, and to demonstrate its behaviour under injected faults
    rather than only under nominal conditions;
  \item to characterise the reachable workspace of the mounted configuration,
    and to determine which manipulation tasks it can actually support;
  \item to establish what performance survives when the master arm's sensing
    degrades, since a substantial fraction of its channels are unreliable;
  \item to design an experimental protocol capable of separating the
    contribution of machine autonomy from that of the human operator.
\end{enumerate}

\section{Challenges}

Three difficulties shaped the work more than any others.

The first is that \emph{the base moves}. Conventional manipulator control
assumes an inertial base; here the base is a standing person. This is a
recognised problem in the supernumerary-limb literature
\todo{CITE: Zhang et al., Motion-Compensation Control of Supernumerary Robotic
Arms Subject to Human-Induced Disturbances, Advanced Intelligent Systems 2024}
but it is normally treated in the single-person case, where the wearer is also
the operator and their own motor commands are available as feed-forward. In
the two-person case that information does not exist.

The second is that \emph{the input is degraded}. The master arm's
potentiometer channels are, in their present physical state, substantially
unreliable, and the system therefore has to remain usable on a subset of its
intended sensing rather than assuming all of it.

The third is that \emph{almost every failure in this system is silent}. An arm
that has stopped moving looks the same whether it is blocked for a good
reason, blocked for a bad one, or not receiving data at all. A great deal of
the engineering effort described in this thesis went into making the system's
internal state legible, and a recurring theme of the results is that
measurement instruments failed more often, and more quietly, than the system
being measured.

\section{Contributions}

The contributions claimed here are limited to those the implementation and its
recorded artefacts can support.

\begin{enumerate}
  \item \textbf{A complete, safety-instrumented dual-arm teleoperation stack}
    for a backpack-mounted pair of Kinova Gen3 arms, comprising master-arm
    sensing, a spherical pose mapping, collision-aware inverse kinematics with
    redundancy resolution, a clearance floor referenced to the wearer's body,
    and a multi-path emergency stop. The behaviour of the safety layer was
    exercised by systematic fault injection rather than by inspection.

  \item \textbf{A workspace characterisation of the mounted configuration},
    including the negative result that the two arms share no reachable
    workspace, which blocks inter-arm handover for geometric reasons that no
    control or planning change can address.

  \item \textbf{A quantified account of degraded-input teleoperation}: which
    capabilities survive the loss of individual sensing channels, and what
    each loss costs.

  \item \textbf{A verification methodology}, with a set of recorded runs
    covering every task, scenario and autonomy condition, checked
    automatically for object handling, gripper actuation and wearer clearance.

  \item \textbf{An experimental design for the two-person configuration},
    including measures for the wearer as a participant rather than as
    apparatus.
\end{enumerate}

Chapter~\ref{ch:results} states explicitly, for each result, whether it was
verified in simulation, exercised only against mock hardware, or validated on
physical arms. \todo{CONFIRM: that the contribution list matches what the
final Results chapter supports, once the outstanding measurements are taken}
